%% Esempio per lo stile supsi
% \documentclass[twoside]{supsistudent} 
% tolto sennò iniziava una pagina bianca dopo il frontespizio

\documentclass[oneside]{supsistudent} 
\let\cleardoublepage\clearpage



% per settare noindent
\setlength{\parindent}{0pt}


% Crea un capitolo senza numerazione che pero` appare nell'indice %
\newcommand{\problemchapter}[1]{%
  \chapter*{#1}%
  \addcontentsline{toc}{chapter}{#1}%
\markboth{#1}{#1}
}

% Numerazione delle appendici secondo norma
\addto\appendix{
\renewcommand{\thesection}{\Alph{chapter}.\arabic{section}}
\renewcommand{\thesubsection}{\thesection.\arabic{subsection}}}

\setcounter{secnumdepth}{5} 	%per avere più livelli nei titoli
\setcounter{tocdepth}{5}		%per avere più livelli nell'indice


\titolo{Magic Eraser}
\studente{Abate Luca \vspace{1em}\\ Vavassori Enrico}
\relatore{Gremlich Giuliano}
\correlatore{Quattrini Andrea}
\committente{PLACEHOLDER}
\corso{Ingegneria Informatica TP}
\modulo{Progetto di Semestre}
\anno{2026}



\begin{document}

\pagenumbering{alph}
\maketitle
\onehalfspacing
\frontmatter


\pagenumbering{roman}
\tableofcontents
% \listoffigures					% Opzionale
% \listoftables					% Opzionale

\newpage
\mainmatter
\pagenumbering{arabic}
\setcounter{page}{1}


%-----------------1-----------------

\chapter{Introduzione}

Magic Eraser is a web-based AI image 
editing platform that allows users to manage projects, upload 
images, and apply generative editing workflows through an 
interactive laboratory interface. Instead of focusing only on 
the removal of unwanted objects or people, the system 
supports a broader pipeline based on AI models, 
including prompt-based segmentation, mask-based removal, 
and generative filling. Users can identify specific areas 
of an image, remove selected content, and regenerate missing 
regions in a visually coherent way, while preserving texture, 
lighting, and overall context. The platform also stores editing 
pipelines and results, enabling users to organize, revisit, and 
continue their image editing work in a structured and reproducible 
manner.
\\The idea of Magic Eraser is to be a research platform to
analyze in an easier way the behavior of generative models.




%-----------------2-----------------

\chapter{Obiettivi e Requisiti}
RIASSUNTO bla bla bla

\section{Obiettivo generale}
\lipsum[11]

\section{Obiettivi specifici}
\lipsum[11]

\section{Requisiti funzionali}
\lipsum[11]

\section{Requisiti non funzionali}
\lipsum[11]


%------------------3-----------------

\chapter{State of the Art e Analisi Bibliografica}
RIASSUNTO bla bla bla

\section{Image inpainting}
\lipsum[11]

\section{Segmentazione delle immagini}
\lipsum[11]

\subsection{Ruolo della mask}
\lipsum[11]

\subsection{ConvexHull}
\lipsum[11]

\section{Modelli generativi per l'editing}
\lipsum[11]

\section{Soluzioni esistenti}
\lipsum[11]

\section{Posizionamento del progetto}
\lipsum[11]


%------------------4-----------------

\chapter{Analisi della Situazione As-Is}
RIASSUNTO bla bla bla

\section{Workflow manuale tradizionale}
\lipsum[11]

\section{Problemi e limiti}
\lipsum[11]

\section{Bisogni dell'utente target}
\lipsum[11]


%------------------5-----------------

\chapter{Analisi e Progettazione della Soluzione To-Be}
RIASSUNTO bla bla bla

\section{Visione generale della soluzione}
\lipsum[11]

\section{Architettura software}
\lipsum[11]

\subsection{Frontend}
\lipsum[11]

\subsection{Backend}
\lipsum[11]

\subsection{Database}
\lipsum[11]

\subsection{Provider AI}
\lipsum[11]

\subsection{Storage immagini}
\lipsum[11]

\section{Modello dei dati}
\lipsum[11]

\subsection{Profile}
\lipsum[11]

\subsection{Project}
\lipsum[11]

\subsection{Image}
\lipsum[11]

\subsection{Process}
\lipsum[11]

\subsection{LaboratoryPipeline}
\lipsum[11]

\subsection{LaboratoryPipelineStep}
\lipsum[11]

\section{Flusso delle pipeline}
\lipsum[11]

\section{Gestione dei progetti}
\lipsum[11]

\section{Gestione dei processi}
\lipsum[11]


%------------------6-----------------

\chapter{Metodologia e Scelte Tecnologiche}
RIASSUNTO bla bla bla

\section{Metodo di sviluppo}
\lipsum[11]

\section{Scelte frontend}
\lipsum[11]

\section{Scelte backend}
\lipsum[11]

\section{Scelte AI}
\lipsum[11]

\subsection{Segmentazione da prompt}
\lipsum[11]

\subsection{Rimozione con maschera}
\lipsum[11]

\subsection{Generazione da prompt}
\lipsum[11]

\section{Persistenza e deployment locale}
\lipsum[11]


%------------------7-----------------

\chapter{Implementazione}
RIASSUNTO bla bla bla

\section{Struttura del progetto}
\lipsum[11]

\subsection{frontend/}
\lipsum[11]

\subsection{backend/}
\lipsum[11]

\subsection{backend/app/api/routers/}
\lipsum[11]

\subsection{backend/app/services/}
\lipsum[11]

\subsection{backend/app/repositories/}
\lipsum[11]

\subsection{backend/app/processes/}
\lipsum[11]

\subsection{backend/app/models\_registry/}
\lipsum[11]

\section{Frontend}
\lipsum[11]

\subsection{Login e signup}
\lipsum[11]

\subsection{Home}
\lipsum[11]

\subsection{Gallery}
\lipsum[11]

\subsection{Laboratory}
\lipsum[11]

\subsection{Pipelines}
\lipsum[11]

\subsection{Profile}
\lipsum[11]

\section{Backend e API}
\lipsum[11]

\subsection{Images}
\lipsum[11]

\subsection{Projects}
\lipsum[11]

\subsection{Profile}
\lipsum[11]

\subsection{Processes}
\lipsum[11]

\subsection{Laboratory pipelines}
\lipsum[11]

\section{Gestione immagini}
\lipsum[11]

\section{Sistema dei processi AI}
\lipsum[11]

\subsection{segment\_from\_prompt}
\lipsum[11]

\subsection{remove\_with\_mask}
\lipsum[11]

\subsection{generate\_from\_prompt}
\lipsum[11]

\subsection{Registry dei processi}
\lipsum[11]

\subsection{Validazione input}
\lipsum[11]

\subsection{Scelta modello}
\lipsum[11]

\subsection{Adapter verso provider AI}
\lipsum[11]

\section{Pipeline di laboratorio}
\lipsum[11]


%------------------8-----------------

\chapter{Verifica, Test e Validazione}
RIASSUNTO bla bla bla

\section{Strategia di test}
\lipsum[11]

\section{Test automatici}
\lipsum[11]

\subsection{Test API progetti}
\lipsum[11]

\subsection{Test API profilo}
\lipsum[11]

\subsection{Test pipeline}
\lipsum[11]

\subsection{Test componenti frontend}
\lipsum[11]

\subsection{Test utility}
\lipsum[11]

\section{Validazione funzionale}
\lipsum[11]

\section{Valutazione qualitativa dei risultati AI}
\lipsum[11]


%------------------9-----------------

\chapter{Analisi Critica dei Risultati}
RIASSUNTO bla bla bla

\section{Valutazione rispetto agli obiettivi}
\lipsum[11]

\section{Qualita dei risultati visivi}
\lipsum[11]

\subsection{Oggetti piccoli}
\lipsum[11]

\subsection{Sfondi uniformi}
\lipsum[11]

\subsection{Sfondi complessi}
\lipsum[11]

\subsection{Oggetti parzialmente occlusi}
\lipsum[11]

\subsection{Prompt ambigui}
\lipsum[11]

\section{Limiti tecnici}
\lipsum[11]

\subsection{Latenza}
\lipsum[11]

\subsection{Costi API}
\lipsum[11]

\subsection{Affidabilita del provider}
\lipsum[11]

\subsection{Assenza di elaborazione locale}
\lipsum[11]

\subsection{Gestione errori}
\lipsum[11]

\subsection{Scalabilita}
\lipsum[11]

\section{Benefici apportati}
\lipsum[11]


%------------------10-----------------

\chapter{Conclusioni e Sviluppi Futuri}
RIASSUNTO bla bla bla

\section{Riassunto del lavoro svolto}
\lipsum[11]

\section{Valutazione finale}
\lipsum[11]

\section{Competenze acquisite}
\lipsum[11]

\section{Possibili generalizzazioni}
\lipsum[11]

\section{Sviluppi futuri}
\lipsum[11]

\subsection{Editor manuale delle maschere}
\lipsum[11]

\subsection{Confronto tra piu modelli}
\lipsum[11]

\subsection{Undo e redo}
\lipsum[11]

\subsection{Esportazione delle pipeline}
\lipsum[11]

\subsection{Metriche quantitative sulla qualita}
\lipsum[11]

\subsection{Elaborazione batch}
\lipsum[11]

\subsection{Utilizzo di modelli locali}
\lipsum[11]

\subsection{Collaborazione multiutente}
\lipsum[11]

\subsection{Versionamento dei risultati}
\lipsum[11]

\subsection{Dashboard per costi e tempi di elaborazione}
\lipsum[11]


%------------------Bibliografia-----------------

\bibliographystyle{unsrt}
\bibliography{bibliografia}


%------------------Appendici-----------------

\appendix

\chapter{Manuale Utente}
RIASSUNTO bla bla bla

\section{Login}
\lipsum[11]

\section{Creare progetto}
\lipsum[11]

\section{Caricare immagine}
\lipsum[11]

\section{Usare laboratorio}
\lipsum[11]

\section{Salvare pipeline}
\lipsum[11]

\chapter{API principali}
RIASSUNTO bla bla bla

\section{Elenco endpoint}
\lipsum[11]

\chapter{Configurazione e Avvio}
RIASSUNTO bla bla bla

\section{Docker Compose}
\lipsum[11]

\section{Variabili ambiente}
\lipsum[11]

\section{Avvio frontend e backend}
\lipsum[11]

\chapter{Test}
RIASSUNTO bla bla bla

\section{Output test}
\lipsum[11]

\section{Copertura}
\lipsum[11]

\section{Elenco casi}
\lipsum[11]

\end{document}
